\section{Introduction}\label{sec:intro}

\subsection{The experiment and the folklore}
Take a long strip of paper, tie a simple overhand knot in it, pull the two ends gently, and flatten. The knot settles into a shape that looks, strikingly, like a regular pentagon. The observation is old recreational mathematics; the usual explanation exhibits a one-parameter family of flat folded configurations and notes that the regular pentagon sits at its contact-limited endpoint, and the construction has been formalized algebraically in computational origami by Ida, Ghourabi, and Takahashi \cite{ida2015}. Kauffman \cite{kauffman2004} turned the folklore into a precise question. Modelling the flattened knot as a folded ribbon and measuring the \emph{truncation length}---cut the two tails flush with the outer boundary of the flat knot and unfold---he computed, for the pentagonal trefoil,
\begin{equation}\label{eq:kauffman-ratio}
\frac{L}{W}=\frac{4}{\sqrt{5-2\sqrt5}}=5.5055276819\ldots,
\end{equation}
and conjectured this value to be the minimum for the flattened truncated trefoil. In the two decades since, the conjecture has attracted upper-bound constructions in related models but, to our knowledge, no optimality result of any kind: no criticality, no local minimality.

\subsection{Two ways to pull tight}
The folklore sentence---\emph{pull it tight and you get a pentagon}---hides an ambiguity which turns out to be the heart of the matter. Two quantities have equal claim to the name \emph{tightness}:
the truncation length $\trunc$, the total material needed to \emph{build} the flat knot with tails cut flush, a functional of manufacture; and the core length $\core$, the material trapped between the first and the last crease, which is what pulling on the tails monotonically extracts, a functional of tightening. On the standard symmetric family the two narratives coincide, and both end at the pentagon. On the full configuration space they do not. This paper shows that the two functionals select \emph{different} strict local optima on the branch, and that the folklore statement is precisely true for one of them and precisely false for the other.

\subsection{Main results}
We develop the kinematic theory from the bottom up and push every claim to the strongest verification grade available to us.

Within the stated model---piecewise $C^2$ width-one strips, complete movable creases, zero thickness---the tying combinatorics is forced. A tied flat folded long strip needs at least three creases, and with exactly three the knot type is the trefoil and the entire tying data---which piece pairs cross, in which material order, with which over/under pattern, with which crease-slope signs, under which layer order---is unique up to mirror image, by a self-contained argument citing no knot tables. The configuration space of the resulting branch is a five-dimensional semialgebraic set whose non-penetration boundary we classify completely: six endpoint-sliver walls together with a disjunctive taco pair, the latter revealed to be corner contacts of crease pairs.

The pentagon is a contact-mode junction---informally, a mountain pass---for pulling, and the certified tight state beyond it is algebraically wild. On the symmetric slice, tightening terminates at the regular pentagon---the folklore, proved, and machine-checked in Lean---but in the full five-dimensional space an exact Farkas computation exhibits a legal first-order descent direction, and descent continues along a second symmetric family---feasibly and strictly monotonically all the way, by exact resultant--Sturm certificates. That family reduces to a plane curve $N_5(x,h)=0$ whose unique critical point $\td$ generates a coordinate field of degree twelve over $\Q$, with defining sextic of Galois group $S_6$: \emph{the state reached by the certified pulling path cannot be written in radicals, its $x$-coordinate already having none.} Existence, identification, KKT blocking with strictly positive multipliers, and strict local minimality via an exact reduced Hessian are all certified over $\Q(x^\dagger)$.

Kauffman was right, locally, and no other blocked state is certified. At the pentagon the truncation length has exact gradient $(0,\varphi^3,-\varphi^2,\varphi,0)$ in powers of the golden ratio $\varphi=(1+\sqrt5)/2$, with exact Farkas certificates on an invertible wall system in both taco branches, so the pentagon is a \emph{strict} local minimizer of $\trunc$ on the three-crease branch, with linear growth; since $\trunc(\td)=5.582\ldots>5.506\ldots$, the new state is no counterexample to Kauffman's conjecture. Two further theorems close the picture: for open states of the three-crease branch the sheet-separation layer axioms are equivalent to realizability by embedded zero-thickness strips in slabs $\R^2\times[0,\varepsilon]$ with $\varepsilon\to0$, the zero-clearance case being certified at the pentagon and at $\td$, and a stratified census of all $215$ active-set strata---run through two independent implementations and finished with per-candidate interval certificates---finds exactly one certified blocked state, namely $\td$, with the residual components localized and graded.

The broadest message is the dichotomy: \emph{tightest is not a property of a shape but of a functional.} One wall system, one moduli space, and a single change of objective move a strict local optimum from the golden-ratio world of the pentagon to an $S_6$-algebraic notched heptagon. Precise statements are Theorems A--H in Section~\ref{sec:mainresults}.

The program is not confined to the trefoil. At four creases the same enumeration machinery yields a unique irreducible four-crossing code realizing the figure-eight knot, and on its branch the pulling functional is again blocked at an exactly solvable state---this time in the quadratic field $\Q(\sqrt{15})$, with core length $4\sqrt{15}/3$ and a centrally symmetric hexagonal outline tiled by the core passes without a notch; its truncation length recovers exactly Kauffman's conjectured figure-eight value $32/\sqrt{15}$ \cite{kauffman2004}---the value is his, the exact realization, optimality, and reachability are new. The contrast is instructive: one more crossing, and the certified blocked local minimizer---in its numerically identified cell---drops from degree twelve and unsolvable to degree two. Section~\ref{sec:figure8} develops this second instance end to end.

\subsection{Predictions}
The dichotomy is falsifiable. The certified terminal state of the pulling path differs from the pentagon by a V-notch on one edge, with mouth $0.0845w$, legs $0.0514w$, depth $0.0293w$, and reflex angle $249.52^\circ$ for a strip of width $w$---all four determined by exact elements of the sextic subfield $\Q((x^\dagger)^2)$, namely the three squared lengths and the cosine of the angle. On a $2\,$cm tape this is a notch about $1.7\,$mm wide and $0.6\,$mm deep. Moreover, along one certified opening direction our constant-height hairpin construction realizes nearby open states at cost $\Delta\ell\lesssim3.5\,\varepsilon$, the coefficient combining the exact multiplier sum $\sum_j\lambda_{w_j}=4.11\ldots$ with the slab-budget rate; the KKT multipliers are the unit prices of the ramp clearances. This is a construction-specific $O(\varepsilon)$ upper bound; no universal lower bound, and no fixed-thickness exclusion of $\td$, is proved. Real paper may instead lock at the pentagon, through friction at eight simultaneous contacts and pinned creases; that would itself be this model's explanation of the folklore, and the notch remains the discriminating signal rather than a guaranteed outcome. See Section~\ref{sec:shape}.

\subsection{Methods and verification}
Every structural claim in this project was probed numerically before being asserted, then promoted along a fixed ladder: multistart numerics; high-precision identification by PSLQ \cite{fba1999} over the relevant number field; exact verification in quotient fields $\Q[x]/(p)$; interval arithmetic (interval Newton, Krawczyk) for signs and existence; and, for the algebraic core, machine-checked proofs in Lean~4/Mathlib \cite{mathlib2020}---a discipline that at one point overturned one of our own intermediate claims and thereby uncovered $\td$ (Remark~\ref{rem:methods}). Internal verification is not journal peer review; the complete suite, with pinned dependency versions and reproduction instructions, accompanies the paper as ancillary material \cite{anc}, and Appendix~\ref{app:suite} describes it.

\subsection{Related work}
Kauffman's paper \cite{kauffman2004} is the origin of the truncated functional; among the works surveyed here we have found no prior optimality or non-optimality result for the pentagon. Ida, Ghourabi, and Takahashi \cite{ida2015} formalized the classical three-fold construction, proving algebraically---with Gr\"obner-basis certificates---that the resulting polygon is regular when the terminal contact points coincide, and treating pentagonal and heptagonal knot folds; their work concerns the construction and verification of a prescribed fold, and does not formulate the deformation space, a tightening functional, or the contact-constrained optimality problem studied here. The modern folded-ribbon literature studies \emph{closed} folded ribbonlength through upper-bound constructions across all diagrams---Denne and coauthors \cite{denne2020}, with the survey \cite{denne2018}; recently Chen, Denne, Patterson, and Patterson \cite{chen2025}; general linear upper bounds in the crossing number by Kim, No, and Yoo \cite{kimnoyoo2026}, subsequently sharpened through Yoo's four-page index \cite{yoo2026b}; bounded ribbonlength for knot families and multi-twist M\"obius bands by Denne and Patterson \cite{dennepatterson2025}; and, for alternating links with bipartite dual graphs, Yoo \cite{yoo2026}---a different quantity and a different problem from our fixed-branch, open and truncated kinematics. The immersed thread \cite{immersed2020,ayala2025} proves exact minima for $C^1$ immersed disk diagrams via vertex tension balance, the immersed cousin of our wall-KKT balance, with no folds and hence no overhand branch. Layer conditions for flat folding go back to Justin \cite{justin1997}, Hull \cite{hull2002}, and Demaine--O'Rourke \cite{demaine2007}; monotone-path formalizations to Connelly, Demaine, and Rote \cite{cdr2003}; contact-force criticality for ropelength to Schuricht and von der Mosel \cite{svdm2004} and Cantarella, Kusner, and Sullivan \cite{cks2002}. The combination pursued here---a fixed tying branch, a moduli space with classified contact walls, two tightness functionals, and end-to-end certificates---appears to be new.

\subsection{Organization}
Section~\ref{sec:mainresults} states the main theorems. Sections~\ref{sec:model}--\ref{sec:census} develop the model, the walls, the tightening theory, the second family, the dichotomy, realization, the shape, and the census. Section~\ref{sec:figure8} carries the program to the figure-eight knot. Section~\ref{sec:open} collects open problems. Appendix~\ref{app:suite} documents the verification suite and Appendix~\ref{app:lean} the Lean formalization.

\section{Main theorems}\label{sec:mainresults}

Throughout, the material is a width-one long strip folded flat at three complete creases along the overhand (long trefoil) branch $\mathcal B$. States are pairs $(f_\theta,\lambda)$ of a planar placement and a layer order, with passes labelled $A,B,C,D$ in material order and coordinates $\theta=(x,g,h,a,b)\in\M\subset\R^5$. On the computational chamber $\M=\Mhat$ the working constraints are the six sliver walls $w_1,w_4,w_5,w_6,w_7,w_8$ and the unwindowed taco pair $w_9,w_{10}$, feasibility being $\{w_9\ge0\}\cup\{w_{10}\ge0\}$; on the honest moduli the taco faces are windowed, as stated in Theorem~\ref{mainthm:moduli}; the layer order $B>D>A>C$ is forced by the branch data. Creases are movable fold lines: a tightening path varies crease positions and angles continuously while preserving the branch data; there are no material-pinned creases and no refolding. A valid state is \emph{blocked} when $\nabla\core$ lies in the cone generated by the gradients of its active walls, taken within the relevant taco branch. Precise definitions occupy Section~\ref{sec:model}.

\begin{mainthm}[fold number and rigidity of the tying]\label{mainthm:fold}
A tied flat folded long strip, in the general-position class of Definition~\textup{\ref{def:gp}}, has at least three creases. With exactly three, the knot type is the left- or right-handed trefoil, and the full tying data---crossing pairs, material passage word $\mathtt{XYZXYZ}$, alternating over/under pattern, alternating crease-slope signs $\epsilon=(+,-,+)$, layer order $B>D>A>C$---is unique up to mirror image. The proof cites no knot classification tables.
\end{mainthm}

\begin{mainthm}[moduli and walls]\label{mainthm:moduli}
The honest moduli $\Mhon$, its positive-taco locus $\Mplus$, and the computational chamber $\M=\Mhat$ are five-dimensional semialgebraic sets with $\Mplus\subseteq\Mhat\subseteq\Mhon$. The non-penetration boundary of $\Mhon$ consists of the six sliver walls and the two taco faces windowed by $0\le u\le1$ and $0\le t\le1$; the chamber is cut out by their unwindowed extensions, capture faces that bound no contact; all eight wall functions admit polynomial representatives with integer coefficients, and $\core$ is a rational function of $\theta$. The taco walls are corner contacts of crease pairs: $w_9=w_{10}=0$ exactly when $Q_0=F$, a three-dimensional manifold $\mathcal F$ whose intersection with $\{w_1=w_4=0\}$ is the standard slice, on which all six of $w_1,w_4,w_7,w_8,w_9,w_{10}$ vanish identically.
\end{mainthm}

\begin{mainthm}[the pentagonal contact-mode junction]\label{mainthm:pass}
On the symmetric slice, tightening terminates at the regular pentagon. At the pentagon, $\nabla\core=\tfrac12\mathbf 1$ exactly, and $\nabla\core$ lies in no subcone of the active wall gradients in either taco branch: an explicit first-order descent direction exists along a second symmetric family, and the pentagon is the crossing corner of two symmetric one-parameter families. Along the second family, parametrized by the branch through $\td$ of the plane curve $N_5(x,h)=0$ of Theorem~\textup{\ref{mainthm:dagger}}, the four remaining walls satisfy $w_7=w_8>0$ and $w_9=w_{10}>0$ strictly between the endpoints, vanishing exactly at the pentagon, and $\core$ is strictly monotone with derivative vanishing exactly at $\td$: exact endpoint identities, five root-free resultant--Sturm certificates on an interval containing $[x^\dagger,x_0]$, and an interval-certified interior sample. Hence descent from the pentagon along the second family reaches $\td$ through valid states of the branch---branch preservation being certified by a further family of root-free resultants---with strictly decreasing core length, and the reverse of the opening segment at the pentagon supplies the slack start: $\td$ is reachable from slack states, and \emph{tight} in the sense of Definition~\textup{\ref{def:pull}}.
\end{mainthm}

\begin{mainthm}[the tight state]\label{mainthm:dagger}
The second family reduces to the plane curve $N_5(x,h)=0$ with $a=h$, $b=(x^2+1)/(2x)$ and $\core=h+b$; write $\ell(x)$ for the resulting restriction of $\core$ to the family. Its unique critical point $\td$ satisfies: $x^\dagger$ has minimal polynomial $p_x=x^{12}+4x^{10}-7x^4+20x^2-2$; the field $\Q(\td)=\Q(x^\dagger)$ has degree twelve; the Galois group of the even subfield generator is $S_6$, so $\td$ admits no radical expression; and $\sqrt5\notin\Q(x^\dagger)$. The KKT multipliers at $\td$ are explicit elements of $\Q(x^\dagger)$, strictly positive, with $\lambda_{w_1}=\lambda_{w_4}$ identically and $\lambda_{w_5}$ in the sextic subfield; and $\td$ is a strict local minimizer of $\core$, the reduced Hessian on the one-dimensional critical cone being the exact element
\begin{multline*}
\ell''(x^\dagger)=\Bigl(\tfrac{2483}{24}x^{11}+\tfrac{849}{2}x^{9}+\tfrac{175}{4}x^{7}+\tfrac{53}{12}x^{5}-\tfrac{17377}{24}x^{3}\\
+\tfrac{7977}{4}x\Bigr)\Big|_{x=x^\dagger}=618.248\ldots>0.
\end{multline*}
\end{mainthm}

\begin{mainthm}[first optimality toward Kauffman's conjecture]\label{mainthm:kauffman}
At the pentagon, $\nabla\trunc=(0,\varphi^3,-\varphi^2,\varphi,0)$ exactly, and both taco branches admit exact Farkas certificates with five strictly positive multipliers on an invertible five-by-five wall system. Hence every nonzero feasible direction strictly increases $\trunc$, and the pentagon is a strict local minimizer of the truncation length on the three-crease branch, with linear growth. In contrast, $\trunc(\td)=5.5818\ldots>4/\sqrt{5-2\sqrt5}$: the two functionals select different strict local optima.
\end{mainthm}

\begin{mainthm}[realization completeness for the three-crease branch]\label{mainthm:realization}
For $\mathcal B$-states, in the sheet-separation layer model: necessity is unconditional---every limit realization by embeddings into $\R^2\times[0,\varepsilon_k]$ with $\varepsilon_k\downarrow0$ satisfies the combinatorial axioms; open states satisfy the converse unconditionally, admitting strong $\varepsilon$-realizations for all $\varepsilon<\varepsilon_0(\theta)=\frac{2}{\pi\Lambda}\min(m_{\mathrm{mat}},2c_{\mathrm{nb}},c_{\mathrm{cr}})$; for zero-clearance states the converse holds under the opening hypothesis of the text, which is exactly certified at the pentagonal and dagger states, so that both are limits of open realized states, with exact rank and positivity certificates.
\end{mainthm}

\begin{mainthm}[stratified census]\label{mainthm:census}
\textup{(1) [proved]} Every tight state satisfying the Mangasarian--Fromovitz constraint qualification is blocked, and every blocked point appears as a KKT root of one of $215$ active-set strata---all active sets of at most five walls, both-taco sets included as redundant super-captures, less the four subsets of $\{w_4,w_6\}$ excluded by $x$-independence---and by Carath\'eodory at most one taco wall is ever needed for capture. \textup{(2) [exact and interval certificates]} Within these strata, the pentagon appears as a root in $20$ of them and is excluded wholesale by the single exact Farkas certificate of Theorem~\textup{\ref{mainthm:pass}}; all $72$ candidates admitting certifiable regular roots are excluded by interval certificates; $\td$ is the unique \emph{certified} blocked state. 
\end{mainthm}

\begin{compres}[Computational census evidence]\label{compres:censusintro}
That no blocked state escapes the remaining numerical components---per-stratum root coverage of the double-implementation harvest, absence of valid non-MFCQ points, on-branch scale calibration, and the $32$ Krawczyk-inconclusive candidates, each of which violates a true wall numerically---is computational evidence, localized in Section~\ref{sec:census}. Reachability of $\td$ by a general descent from an arbitrary slack state, as opposed to along the certified path of Theorem~\textup{\ref{mainthm:pass}}, is a separate, numerically supported statement.
\end{compres}

\begin{mainthm}[an exact blocked hexagon on the figure-eight branch]\label{mainthm:hex}
At four creases the enumeration of Theorem~\textup{\ref{mainthm:fold}} yields a unique irreducible four-crossing code, and exactly two of its sixteen over/under assignments---a mirror pair---realize the figure-eight knot $4_1$. In the numerically identified cell of the resulting seven-dimensional branch described below---an $8\times10^6$-sample sweep finds valid width-one flat folded states, all in a single cell, self-mirror under material reversal, with layer order $P_1>P_4>P_2>P_0>P_3$ up to the top-bottom flip---the following centrally symmetric state $\theta^*_{4_1}$, exactly solvable in $\Q(\sqrt{15})$, is a valid, reachable, strict local minimizer of $\core$:
\begin{gather*}
\cos\tau_1=-\tfrac14,\qquad \tau_2=2\pi-2\tau_1=-\tau_3,\qquad \tau_4=-\tau_1,\\
(L_1,L_2,L_3)=\tfrac{1}{\sqrt{15}}\,(6,8,6),\qquad \core=\tfrac{4\sqrt{15}}{3}.
\end{gather*}
Its eight active contacts vanish exactly, its KKT multipliers on an invertible seven-generator subsystem of the eight active gradients are positive elements of $\Q(\sqrt{15})$, and first-order sufficiency upgrades blocking to \emph{strict} local minimality of $\core$. Its core outline is a centrally symmetric hexagon with sides $\tfrac{2\sqrt{10}}5$ and $\tfrac{4}{\sqrt{15}}$ and interior angles $\arccos(-\tfrac14)$ and $\pi-\tfrac12\arccos(-\tfrac14)$, which the three core passes tile without a notch, exactly; and the truncation length is $\trunc(\theta^*_{4_1})=\tfrac{32\sqrt{15}}{15}$, each truncation stub equal to $L_1$, so that $\trunc/\core=\tfrac85$ exactly; the value $32/\sqrt{15}$ is precisely the figure-eight ribbon length conjectured by Kauffman \cite{kauffman2004}. The state itself is certified valid: all thirty-nine candidate contact margins of the branch inventory, evaluated at $\theta^*_{4_1}$ in $\Q(\sqrt{15})$, come out exactly zero for the eight active contacts and strictly positive for the other thirty-one, so the existence of valid states is exact, witnessed by $\theta^*_{4_1}$. Moreover the exact opening direction furnishes a slack start whose reversal reaches $\theta^*_{4_1}$, which is therefore reachable and \emph{tight} in the sense of Definition~\textup{\ref{def:pull}}. Exhaustiveness of the cell within the sweep and convergence of \emph{generic} descent are numerical.
\end{mainthm}
